% !TeX root = ../../Book1.tex
\section{导数测度的分解}
\label{b1:sec:2803}

分布导数把一个函数的全部一阶变化保存在同一向量测度中。
但这些变化可以分布在体积中，也可以集中在跳跃面上，还可能既不占体积、又不由跳跃面承载。
本节区分这三种情形。前两种分别连接 Sobolev 梯度与有限周长集合，
第三种则解释连续函数为何仍可能没有可积的弱导数。

以下 \(\Omega\subseteq\mathbb R^d\) 为开集，\(u\in BV(\Omega)\) 为实值。
跳跃值与法向沿用\secref{b1:sec:2802}的约定：
\(u^+>u^-\)，\(\nu_u\) 从低值侧指向高值侧。
深层结构的准确陈述可参看 Alberti、Mantegazza 的
《A Note on the Theory of SBV Functions》\cite[Remark 1.2]{AlbertiMantegazza1997SBV}；
系统学习可继续阅读 Ambrosio、Fusco、Pallara 的
《Functions of Bounded Variation and Free Discontinuity Problems》%
\cite[Chapter 3]{Ambrosio2000Functions}。
本节把测度分解的直接后果完整推出，而将近似微分的识别和链律中的深层步骤明确列为证明概要。
% 实读作者公开正式论文 Alberti--Mantegazza 1997：
% 作者版第2--3页（刊印376--377）Remark 1.2；
% 作者版第4--5页（刊印378--379）Remark 2.1；
% 其§2.3--2.5另供下一节紧性使用。AFP只核实出版与章节导航，无全文阅读声明。
% 原文Su表示无近似极限集；本书保留Su与真正两侧跳跃集Ju的区别。

\subsection{绝对连续部分}
\label{b1:subsec:280301}

先不使用函数的几何性质。对 \(Du\) 的每个分量应用 Lebesgue 分解，
得到唯一的表示
\begin{equation}\label{b1:eq:bv-lebesgue-decomposition}
 Du=g\,m_d+D^su,\quad g\in L^1(\Omega;\mathbb R^d),\quad
 |D^su|\perp m_d.
\end{equation}
记 \(D^au=g\,m_d\)，称它为导数测度的绝对连续部分。
具体说，各分量的奇异部分分别集中在某个体积零的 Borel 集上；
取这有限多个集合的并 \(N\)，便有
\[
 D^au=Du\mathbin{\llcorner}(\Omega\setminus N),
 \quad
 D^su=Du\mathbin{\llcorner}N.
\]
这里限制记号只表示在相应 Borel 集上取测度。
向量极分解给出 \(|D^au|=|g|m_d\)；
不同分量的分解由标量唯一性保证相容。因而这一步只需要第6章的测度论，
尚未证明 \(g\) 与逐点的微分有什么联系。
\symidx{\(D^au\)：\(Du\) 的绝对连续部分}
\symidx{\(D^su\)：\(Du\) 的奇异部分}

\begin{theorem}[有界变差函数的近似梯度]
\label{b1:thm:bv-approximate-gradient}
函数 \(u\) 在 \(m_d\)-几乎处处近似可微。
若在近似极限存在处用该极限选取中心值，则其近似微分可写成
\[
 \operatorname{ap}Du(x)\,h=\nabla u(x)\cdot h,
 \quad \nabla u(x)=g(x)
 \quad m_d\text{-几乎处处},
\]
其中 \(g\) 是\eqref{b1:eq:bv-lebesgue-decomposition}中的密度。
特别地，\(D^au=\nabla u\,m_d\)，且 \(\nabla u\in L^1(\Omega;\mathbb R^d)\)。
\end{theorem}
\begin{proofsketch}
这里不仅需要 Radon--Nikodym 定理，还要证明函数本身在几乎每点有一阶线性近似。
一种路线先用变差的局部平均控制函数振荡，把函数限制在可数个 Lipschitz 片上；
在这些片的密度点应用 Rademacher 定理，再用分布测试识别所得梯度与 \(g\)。
把振荡估计转成覆盖几乎所有点的 Lipschitz 分片，以及最后的密度识别，
是本节未展开的两步，不能由测度有密度直接代替。
上述结论在所引短文 Remark 1.2 和 equation (1.2) 之后明确陈述，
该处将完整论证指向 Evans--Gariepy 原版的 Section 6.1。
\end{proofsketch}

从此对一般有界变差函数写 \(\nabla u\) 时，指这个近似梯度的几乎处处等价类。
它不是 \(Du\) 的另一种写法，也不要求一个任意指定的代表具有经典偏导数。
若 \(u\in W^{1,1}(\Omega)\)，分布导数已经完全由弱梯度表示，
上述唯一性便说明近似梯度与弱梯度几乎处处一致。

\begin{proposition}\label{b1:prop:bv-sobolev-singular-criterion}
对于 \(u\in BV(\Omega)\)，有
\[
 u\in W^{1,1}(\Omega)\quad\Longleftrightarrow\quad D^su=0.
\]
\end{proposition}
\begin{proof}
正向由弱导数的定义和 Lebesgue 分解唯一性得到。
反向若 \(D^su=0\)，则每个分布偏导数都由 \(g_i\in L^1(\Omega)\) 表示。
再加上有界变差函数定义中的 \(u\in L^1(\Omega)\)，恰好满足 \(W^{1,1}\) 的定义。
\end{proof}

\subsection{跳跃部分}
\label{b1:subsec:280302}

\begin{definition}\label{b1:def:bv-jump-derivative}
将导数测度限制到跳跃集 \(J_u\)，所得向量测度记为
\[
 D^ju=Du\mathbin{\llcorner}J_u,
\]
称为导数测度的\term[tiaoyue-bufen]{跳跃部分}[jump part]。
\end{definition}
\symidx{\(D^ju\)：\(Du\) 的跳跃部分}

由\cref{b1:thm:bv-jump-structure}，它具有明确的表面密度：
\begin{equation}\label{b1:eq:bv-jump-density}
 D^ju=(u^+-u^-)\nu_u\,
       \mathcal H^{d-1}\mathbin{\llcorner}J_u.
\end{equation}
跳跃的高度、朝向与承载的面积在这个公式中各占一个位置。
如果改用不排序的两侧记号，同时交换 \(u^+\)、\(u^-\) 并把法向反号，
右侧向量不变。因此测度不依赖人为的定向选择。
由于 \(|\nu_u|=1\)，对每个 Borel 集 \(A\subseteq\Omega\)，
\[
 |D^ju|(A)=\int_{A\cap J_u}(u^+-u^-)\,\mathrm d\mathcal H^{d-1}.
\]
这给出跳跃高度的可积性，但不保证 \(\mathcal H^{d-1}(J_u)<\infty\)：
很多很小的跳跃，可以具有有限总变差和无限总面积。

例如在 \((0,1)\) 内令
\[
 v(x)=\sum_{n=1}^{\infty}2^{-n}\mathbf1_{(2^{-n},1)}(x).
\]
该级数一致绝对收敛，每项的分布导数为 \(2^{-n}\delta_{2^{-n}}\)。
对测试函数逐项积分合法，故
\[
 Dv=\sum_{n=1}^{\infty}2^{-n}\delta_{2^{-n}},
 \quad |Dv|\bigl((0,1)\bigr)=1.
\]
在每个 \(2^{-n}\) 处跳跃高度为 \(2^{-n}\)，其他内部点在一个邻域内函数恒定，
所以 \(J_v=\{\,2^{-n}:n\ge1\,\}\)。
一维的 \(\mathcal H^0\) 是计数测度，因而 \(\mathcal H^0(J_v)=\infty\)。

\subsection{Cantor 部分}
\label{b1:subsec:280303}

跳跃集可数 \((d-1)\)-可整流，因而是体积零集。
所以 \(D^ju\) 是 \(D^su\) 在 \(J_u\) 上的限制，余下的部分仍为奇异测度。

\begin{definition}\label{b1:def:bv-cantor-derivative}
规定
\[
 D^cu=D^su-D^ju
     =D^su\mathbin{\llcorner}(\Omega\setminus J_u),
\]
称 \(D^cu\) 为导数测度的\term[cantor-bufen]{Cantor 部分}[Cantor part]。
\end{definition}
\symidx{\(D^cu\)：\(Du\) 的 Cantor 部分}

\begin{theorem}[有界变差函数的导数测度三分解]
\label{b1:thm:bv-derivative-decomposition}
对于任意实值 \(u\in BV(\Omega)\)，有
\begin{equation}\label{b1:eq:bv-three-part-decomposition}
 Du=\nabla u\,m_d
    +(u^+-u^-)\nu_u\,
      \mathcal H^{d-1}\mathbin{\llcorner}J_u
    +D^cu.
\end{equation}
这三项两两互相奇异。更进一步，若 Borel 集 \(A\subseteq\Omega\)
关于 \(\mathcal H^{d-1}\) 为 \(\sigma\)-有限，则
\[
 |D^cu|(A)=0.
\]
有限近似极限 \(\widetilde u\) 在 \(|D^cu|\)-几乎处处存在。
\end{theorem}
\begin{proofsketch}
分解等式由前面两项的识别和 \(D^cu\) 的定义得到。
若 \(N\) 承载 \(D^su\) 且 \(m_d(N)=0\)，三项分别集中于
\(\Omega\setminus(N\cup J_u)\)、\(J_u\)、\(N\setminus J_u\)，故两两互相奇异。
这里尚不能只由 \(D^cu(J_u)=0\) 推出它不充任何有限面积集。

后一性质是有界变差函数的精细结构所给出的额外结论。
证明须把任意有限 \(\mathcal H^{d-1}\) 集与函数在不同方向上的一维截面比较，
识别其中能承载导数的部分为两侧迹的跳跃部分；
剩余部分在这样的集合上均为零，再由可数并处理 \(\sigma\)-有限集。
这一几何识别没有在本节证明，准确结论见前引 Remark 1.2 的式(1.1)--(1.4)。
结合\cref{b1:thm:bv-jump-structure}，
没有有限近似极限的集合除去 \(J_u\) 只有 \(\mathcal H^{d-1}\) 零测余项，
因此它也是 \(|D^cu|\) 的零测集。这给出最后一句。
\end{proofsketch}

\begin{proposition}\label{b1:prop:bv-decomposition-variation}
三部分分解唯一，并且对每个 Borel 集 \(A\subseteq\Omega\)，
\begin{equation}\label{b1:eq:bv-three-part-variation}
 |Du|(A)=\int_A|\nabla u|\,\mathrm dx
   +\int_{A\cap J_u}(u^+-u^-)\,\mathrm d\mathcal H^{d-1}
   +|D^cu|(A).
\end{equation}
\end{proposition}
\begin{proof}
Lebesgue 分解唯一确定第一项。
在 \(J_u\) 上限制剩余测度，唯一确定第二项；相减便唯一确定第三项。
取上述三个互不交的承载集，在每个承载集上 \(Du\) 等于相应分量，
所以其变差也等于该分量的变差。
由变差的可数可加性相加，再用向量密度的变差公式，即得所述等式。
\end{proof}

Cantor 部分的名称不是说它必须集中在三分 Cantor 集上，
而是说它代表不能由体积密度和跳跃面描述的奇异变化。
在一维，它恰好成为无原子的奇异部分，这一点可以不借高维结构定理直接证明。

\begin{proposition}\label{b1:prop:one-dimensional-bv-decomposition}
设 \(I=(a,b)\) 为有界区间，\(u\in BV(I)\)，\(\mu=Du\)。
则存在常数 \(c\)，使
\[
 \widehat u(x)=c+\mu\bigl((a,x]\bigr)
\]
是 \(u\) 的右连续代表。
其左右极限之差等于 \(\mu(\{x\})\)，且
\[
 D^ju=\sum_{x\in I}\mu(\{x\})\delta_x.
\]
这里非零项至多可数，绝对值之和有限。
测度 \(D^cu\) 则是 \(\mu\) 的奇异部分删去全部原子后的剩余测度；
它的累积分布函数连续。
\end{proposition}
\begin{proof}
函数 \(F(x)=\mu\bigl((a,x]\bigr)\) 有界。由符号测度积分的 Fubini 计算，
对 \(\varphi\in C_c^\infty(I)\)，
\[
 \int_a^b F(x)\varphi'(x)\,\mathrm dx
 =\int_I\left(\int_t^b\varphi'(x)\,\mathrm dx\right)\mathrm d\mu(t)
 =-\int_I\varphi\,\mathrm d\mu.
\]
因此 \(D(u-F)=0\)。
\secref{b1:sec:2004}证明的零弱导数判别说明 \(u-F=c\) 几乎处处。
有限测度的从上、从下连续性分别给出 \(F\) 的右连续性和左极限，
并且 \(F(x)-F(x-)=\mu(\{x\})\)。
非零原子至多可数，因为对每个正整数 \(n\)，质量绝对值至少 \(1/n\) 的原子只有有限个。
对有限原子集求和再取上确界，得到
\(\sum_x|\mu(\{x\})|\le|\mu|(I)\)。

累积分布代表在非原子点连续，在原子点具有不相同的两侧极限；
因此它的近似跳跃集恰为非零原子集。
每个原子的有向跳跃等于其有符号质量，给出 \(D^ju\) 的公式。
绝对连续测度没有原子，故原子全部属于 \(\mu\) 的奇异部分。
删去它们以后所得奇异测度无原子，其累积分布在每点左右连续，因而连续。
\end{proof}

例如，习题\ref{b1:ex:14-cantor-function}中的 \(F_C\) 满足
\[
 DF_C=\rho_C,\quad D^aF_C=0,\quad D^jF_C=0,\quad D^cF_C=\rho_C
 \quad\text{在 }(0,1)\text{ 上}.
\]
第一式由刚才的测试积分得出；后面三式来自 \(\rho_C\) 奇异且无原子。
于是函数
\[
 u(x)=x+\mathbf1_{(1/2,1)}(x)+F_C(x)
\]
在同一区间内同时具有三种非零变化：
\(D^au=m_1\)、\(D^ju=\delta_{1/2}\)、\(D^cu=\rho_C\)，总变差为 \(3\)。
这也表明连续性只能排除跳跃，不能排除 Cantor 部分。

最后记录复合函数所需的精细链律。
下一节会用它分别测试跳跃的存在与梯度的极限。

\begin{theorem}[Volpert 链律的 \(C^1\) 版本]
\label{b1:thm:volpert-chain-rule}
设 \(u\in BV(\Omega)\)，\(\Phi\in C^1(\mathbb R)\)，且 \(\Phi'\) 有界。
若 \(m_d(\Omega)<\infty\) 或 \(\Phi(0)=0\)，则 \(\Phi(u)\in BV(\Omega)\)，并且
\begin{equation}\label{b1:eq:volpert-chain-rule}
 D\bigl(\Phi(u)\bigr)
 =\Phi'(\widetilde u)(D^au+D^cu)
  +\bigl(\Phi(u^+)-\Phi(u^-)\bigr)\nu_u\,
    \mathcal H^{d-1}\mathbin{\llcorner}J_u.
\end{equation}
其中 \(\widetilde u\) 在有限近似极限不存在处可任意补值。
\end{theorem}
\term[volpert-lianlv]{Volpert 链律}[Volpert chain rule]在跳跃处使用两端值之差，
而不是把某个任选的点值代入 \(\Phi'\) 后乘上跳跃高度。
\begin{proofsketch}
一维时先用测度累积代表，把连续变化与原子变化分开：
连续部分满足通常的链式乘法，原子在复合后变成
\(\Phi(u(x+))-\Phi(u(x-))\)。
高维论证用有界变差函数的切片定理将各方向的分布导数还原为这些一维导数，
再重组为向量测度。这里省略的是切片定理及其与三个导数部分的相容性，
不是对公式作形式求导。
所用版本见 Alberti 与 Mantegazza 的《A Note on the Theory of SBV Functions》%
\cite[Remark 2.1, equation (2.1)]{AlbertiMantegazza1997SBV}；
本书按既定方向约定显式写出法向因子。
可积性由 \(|\Phi(t)|\le|\Phi(0)|+\|\Phi'\|_\infty\cdot|t|\) 及给定假设保证，
测度右侧的有限性则由三部分变差公式和 \(\Phi\) 的 Lipschitz 界保证。
\end{proofsketch}
